% ============================================================================
% HoloDraw Project Report — Overleaf Ready
% ============================================================================
\documentclass[12pt, a4paper]{article}

% --- Packages ---
\usepackage[utf8]{inputenc}
\usepackage[T1]{fontenc}
\usepackage{lmodern}
\usepackage[margin=2.5cm]{geometry}
\usepackage{graphicx}
\usepackage{amsmath, amssymb}
\usepackage{booktabs}
\usepackage{hyperref}
\usepackage{enumitem}
\usepackage{float}
\usepackage{caption}
\usepackage{subcaption}
\usepackage{xcolor}
\usepackage{fancyhdr}
\usepackage{titlesec}
\usepackage{parskip}

% --- Hyperlink Styling ---
\hypersetup{
    colorlinks=true,
    linkcolor=blue!60!black,
    citecolor=blue!60!black,
    urlcolor=blue!60!black,
}

% --- Header / Footer ---
\pagestyle{fancy}
\fancyhf{}
\fancyhead[L]{\small HoloDraw — Project Report}
\fancyhead[R]{\small \thepage}
\renewcommand{\headrulewidth}{0.4pt}

% --- Section Styling ---
\titleformat{\section}{\large\bfseries}{\thesection.}{0.5em}{}
\titleformat{\subsection}{\normalsize\bfseries}{\thesubsection.}{0.5em}{}

% ============================================================================
\begin{document}

% ---------- TITLE PAGE ----------
\begin{titlepage}
    \centering
    \vspace*{2cm}

    {\Huge\bfseries HoloDraw}\\[0.5cm]
    {\Large A Hand Gesture--Controlled Drawing Application\\with Deep Learning Edge Detection}\\[2cm]

    \rule{\textwidth}{1pt}\\[0.5cm]

    {\large\bfseries Project Report}\\[0.3cm]
    {\large Deep Learning Module}\\[0.5cm]

    \rule{\textwidth}{1pt}\\[2cm]

    % ---------------------------------------------------------------
    %  TODO: Replace with your actual student names and IDs
    % ---------------------------------------------------------------
    \begin{tabular}{l l}
        \textbf{Student Name}  & \textbf{Student ID} \\[4pt]
        \toprule
        Sahan Perera           & XXXXXXX             \\
        % Add additional team members below
        % Team Member 2       & XXXXXXX             \\
        % Team Member 3       & XXXXXXX             \\
        \bottomrule
    \end{tabular}

    \vfill

    {\large \today}
\end{titlepage}

% ---------- TABLE OF CONTENTS ----------
\tableofcontents
\newpage

% ============================================================================
% 1. ABSTRACT
% ============================================================================
\section{Abstract}

HoloDraw is a real-time, gesture-controlled digital drawing application that merges computer vision with deep learning to create an intuitive, hands-free creative tool.
The system operates in two modes: \textbf{HoloDraw}, which enables freehand sketching and geometric shape placement through hand gestures tracked by MediaPipe, and \textbf{HoloPaint}, which leverages a custom-trained \textbf{U-Net} convolutional neural network to extract clean edge contours from arbitrary images, transforming them into digital colouring-book canvases that users can paint over using gestures or mouse input.
The U-Net model was trained on the \textbf{BSDS500} benchmark dataset using weighted binary cross-entropy loss to address severe class imbalance between edge and non-edge pixels.
Post-processing techniques including morphological gap closing, boundary dilation, and perimeter border injection ensure that extracted outlines form closed regions suitable for flood-fill colouring.
A gesture recognition module classifies six distinct hand poses—draw, navigate, select (pinch), undo (fist), redo (thumb-up), and idle—using geometric heuristics with debounce stabilisation.
Cursor smoothing via exponential moving average filtering further reduces hand-tracking jitter for precise painting.
The system achieves real-time performance on consumer hardware, demonstrating the practical integration of deep learning inference into an interactive desktop application.

% ============================================================================
% 2. PROBLEM STATEMENT
% ============================================================================
\section{Problem Statement}

Traditional digital drawing applications rely exclusively on mouse or stylus input, creating a barrier for users who seek a more natural, hands-free interaction paradigm.
While touchscreen-based drawing has improved accessibility, it still requires physical contact with a device.
The challenge addressed by this project is twofold:

\begin{enumerate}[label=\arabic*.]
    \item \textbf{Gesture-Based Drawing:}
    How can a webcam-based hand tracking system be designed to reliably translate hand gestures into precise drawing actions (freehand strokes, shape placement, tool selection, undo/redo) in real time, despite the inherent noise and ambiguity of markerless hand tracking?

    \item \textbf{AI-Powered Edge Detection for Colouring:}
    How can a deep learning model be trained to extract clean, continuous contour lines from arbitrary photographic images such that the resulting outlines form \emph{closed} regions suitable for digital colouring—analogous to a physical colouring book?
\end{enumerate}

The second challenge is particularly demanding because standard edge detectors (e.g., Canny) produce fragmented, non-contiguous edges that leak when flood-filled.
A learned approach is required to produce semantically meaningful, complete boundaries.

% ============================================================================
% 3. DATASET DESCRIPTION AND PREPROCESSING
% ============================================================================
\section{Dataset Description and Preprocessing}

\subsection{BSDS500 Dataset}

The U-Net edge detection model was trained on the \textbf{Berkeley Segmentation Dataset and Benchmark 500 (BSDS500)}~\cite{bsds500}, a widely used benchmark for contour detection and image segmentation.

\begin{table}[H]
    \centering
    \caption{BSDS500 dataset split summary.}
    \label{tab:dataset}
    \begin{tabular}{lcc}
        \toprule
        \textbf{Split} & \textbf{Images} & \textbf{Purpose} \\
        \midrule
        Training       & 200             & Model weight optimisation \\
        Validation     & 100             & Hyperparameter tuning, early stopping \\
        Test           & 200             & Final performance evaluation \\
        \bottomrule
    \end{tabular}
\end{table}

Each image is accompanied by multiple human-annotated boundary maps (typically 4--9 annotators per image).
The consensus ground truth is constructed by \textbf{averaging} all annotators' edge maps and binarising at a 50\% agreement threshold:
\[
    \text{GT}(x, y) =
    \begin{cases}
        1 & \text{if } \frac{1}{N} \sum_{i=1}^{N} B_i(x, y) > 0.5 \\
        0 & \text{otherwise}
    \end{cases}
\]
where $B_i(x, y)$ is the binary boundary annotation from annotator~$i$ and $N$ is the total number of annotators.

\subsection{Preprocessing and Augmentation}

All images are resized to $288 \times 288$ pixels (slightly larger than the model input) before the following augmentation pipeline is applied during training:

\begin{itemize}[nosep]
    \item \textbf{Random crop} to $256 \times 256$ pixels.
    \item \textbf{Random horizontal flip} (50\% probability).
    \item \textbf{Random vertical flip} (50\% probability).
    \item \textbf{Random rotation} ($\pm 15°$, 50\% probability).
\end{itemize}

All geometric augmentations are applied synchronously to both the RGB image and the corresponding ground-truth edge map to maintain spatial alignment.
Images are normalised using ImageNet statistics ($\mu = [0.485, 0.456, 0.406]$, $\sigma = [0.229, 0.224, 0.225]$).

% ============================================================================
% 4. MODEL ARCHITECTURE AND TRAINING DETAILS
% ============================================================================
\section{Model Architecture and Training Details}

\subsection{U-Net Architecture}

The core deep learning component is a \textbf{U-Net}~\cite{ronneberger2015unet}, a fully convolutional encoder--decoder network with skip connections, adapted for edge detection.

\begin{figure}[H]
    \centering
    \small
    \begin{tabular}{cccc}
        \toprule
        \textbf{Stage} & \textbf{Operation} & \textbf{Output Channels} & \textbf{Spatial Resolution} \\
        \midrule
        \multicolumn{4}{c}{\textit{Encoder (Contracting Path)}} \\
        \midrule
        Enc-1 & ConvBlock & 64   & $256 \times 256$ \\
        Enc-2 & MaxPool $\to$ ConvBlock & 128  & $128 \times 128$ \\
        Enc-3 & MaxPool $\to$ ConvBlock & 256  & $64 \times 64$   \\
        Enc-4 & MaxPool $\to$ ConvBlock & 512  & $32 \times 32$   \\
        \midrule
        \multicolumn{4}{c}{\textit{Bottleneck}} \\
        \midrule
        Bottleneck & MaxPool $\to$ ConvBlock & 1024 & $16 \times 16$ \\
        \midrule
        \multicolumn{4}{c}{\textit{Decoder (Expanding Path)}} \\
        \midrule
        Dec-4 & TransConv $\to$ Cat(skip) $\to$ ConvBlock & 512  & $32 \times 32$   \\
        Dec-3 & TransConv $\to$ Cat(skip) $\to$ ConvBlock & 256  & $64 \times 64$   \\
        Dec-2 & TransConv $\to$ Cat(skip) $\to$ ConvBlock & 128  & $128 \times 128$ \\
        Dec-1 & TransConv $\to$ Cat(skip) $\to$ ConvBlock & 64   & $256 \times 256$ \\
        \midrule
        Output & Conv $1\times1$ $\to$ Sigmoid & 1 & $256 \times 256$ \\
        \bottomrule
    \end{tabular}
    \caption{U-Net architecture summary. Each \texttt{ConvBlock} consists of two $3\times3$ convolutions with BatchNorm and ReLU activation.}
    \label{tab:architecture}
\end{figure}

The model contains approximately \textbf{31 million trainable parameters}.

\subsection{Loss Function}

Edge pixels constitute approximately 5\% of all pixels, creating a severe class imbalance.
Standard binary cross-entropy (BCE) would bias the model towards predicting ``no edge'' everywhere.
We employ \textbf{Weighted BCE with Logits}:
\[
    \mathcal{L} = -\frac{1}{HW} \sum_{i=1}^{HW} \Big[ w_{\text{pos}} \cdot y_i \log(\hat{y}_i) + (1 - y_i) \log(1 - \hat{y}_i) \Big]
\]
where $w_{\text{pos}} = \min\!\left(\frac{1 - r}{r}, 15\right)$ and $r$ is the edge pixel ratio.
The weight is capped at 15 to prevent training instability.

\subsection{Training Configuration}

\begin{table}[H]
    \centering
    \caption{Training hyperparameters.}
    \label{tab:training}
    \begin{tabular}{ll}
        \toprule
        \textbf{Parameter} & \textbf{Value} \\
        \midrule
        Optimiser                  & Adam ($\beta_1=0.9$, $\beta_2=0.999$) \\
        Initial learning rate      & $1 \times 10^{-4}$ \\
        Weight decay               & $1 \times 10^{-5}$ \\
        LR scheduler               & ReduceLROnPlateau (patience=5, factor=0.5) \\
        Batch size                 & 8 \\
        Epochs                     & 50 \\
        Gradient clipping          & Max norm = 1.0 \\
        Input resolution           & $256 \times 256$ pixels \\
        \bottomrule
    \end{tabular}
\end{table}

\subsection{Post-Processing Pipeline}

The raw model output is a probability map in $[0, 1]$.
To produce clean, colourable outlines suitable for the HoloPaint application, the following post-processing steps are applied:

\begin{enumerate}[nosep]
    \item \textbf{Thresholding:} Binarise the probability map at a user-adjustable threshold (default 0.3).
    \item \textbf{Morphological Closing:} Apply an elliptical structuring element to bridge small gaps in contour lines, preventing flood-fill leakage.
    \item \textbf{Dilation:} Slightly thicken edges to ensure solid, continuous boundaries.
    \item \textbf{Border Injection:} Draw a solid 2-pixel black border around the entire canvas perimeter, guaranteeing that edge-adjacent regions cannot leak outside the image bounds.
    \item \textbf{Inversion:} Convert to black-edges-on-white-background format for painting.
\end{enumerate}

\subsection{Hand Gesture Recognition System}

The gesture recognition pipeline operates independently of the U-Net and uses \textbf{MediaPipe HandLandmarker} (a pre-trained deep learning model) to detect 21 3D hand landmarks per frame.
A rule-based classifier then maps finger extension patterns to six gestures:

\begin{table}[H]
    \centering
    \caption{Gesture-to-action mapping.}
    \label{tab:gestures}
    \begin{tabular}{lll}
        \toprule
        \textbf{Gesture} & \textbf{Hand Pose} & \textbf{Action} \\
        \midrule
        DRAW     & Index finger only      & Paint on the canvas \\
        NAVIGATE & Index + middle fingers  & Move cursor without painting \\
        SELECT   & Pinch (thumb--index)    & Tool/colour selection \\
        UNDO     & Fist (all fingers curled) & Undo last action \\
        REDO     & Thumb up only           & Redo last undone action \\
        IDLE     & Open palm / unrecognised & Stop current action \\
        \bottomrule
    \end{tabular}
\end{table}

Gesture stability is ensured through a \textbf{debounce filter} requiring consistent detection over multiple consecutive frames before accepting a gesture change, and an \textbf{exponential moving average (EMA)} cursor smoother that eliminates hand-tracking jitter.

% ============================================================================
% 5. RESULTS AND DISCUSSION
% ============================================================================
\section{Results and Discussion}

\subsection{Quantitative Results}

% ---------------------------------------------------------------
%  TODO: Replace the metric values below with your actual numbers
%        from the evaluation notebook output.
% ---------------------------------------------------------------
\begin{table}[H]
    \centering
    \caption{U-Net performance on the BSDS500 test set (200 images).}
    \label{tab:results}
    \begin{tabular}{lc}
        \toprule
        \textbf{Metric}   & \textbf{Score} \\
        \midrule
        Precision          & 0.XXXX \\
        Recall             & 0.XXXX \\
        F1 Score           & 0.XXXX \\
        Best Val. Loss     & 0.XXXX \\
        \bottomrule
    \end{tabular}
\end{table}

% ---------------------------------------------------------------
%  TODO: Uncomment and update the paths below once you place the
%        generated figures in the report/figures/ directory.
% ---------------------------------------------------------------
% \begin{figure}[H]
%     \centering
%     \includegraphics[width=0.95\textwidth]{figures/evaluation_curves.png}
%     \caption{Training and validation loss curves, and validation F1 score over 50 epochs.}
%     \label{fig:curves}
% \end{figure}

% \begin{figure}[H]
%     \centering
%     \includegraphics[width=0.95\textwidth]{figures/comparison_grid.png}
%     \caption{Visual comparison of original images, ground truth edges, and U-Net predictions on the BSDS500 test set.}
%     \label{fig:comparison}
% \end{figure}

% \begin{figure}[H]
%     \centering
%     \includegraphics[width=0.95\textwidth]{figures/threshold_analysis.png}
%     \caption{Effect of varying the binarisation threshold on edge detection output.}
%     \label{fig:threshold}
% \end{figure}

\subsection{Qualitative Analysis}

The trained U-Net produces significantly cleaner and more semantically meaningful edges compared to the Canny fallback.
Key observations include:

\begin{itemize}[nosep]
    \item The model successfully captures \textbf{object boundaries} while suppressing internal texture edges, producing outlines well-suited for colouring.
    \item At lower thresholds (0.1--0.2), finer details such as facial features and clothing folds are preserved, at the cost of increased noise.
    \item At higher thresholds (0.5--0.7), only the most prominent structural contours remain, yielding a ``clean colouring book'' aesthetic.
    \item The morphological post-processing pipeline effectively closes small gaps in the model output, as evidenced by successful flood-fill operations on random internet images.
\end{itemize}

\subsection{System Performance}

The complete pipeline (hand tracking + gesture recognition + edge detection + rendering) runs at approximately 25--30 FPS on a consumer laptop (Apple M4), demonstrating real-time viability.
The U-Net inference (CPU-only) adds approximately 50--80 ms per image during the initial edge extraction phase but does not impact the painting loop, as edge detection is performed only once per loaded image.

% ============================================================================
% 6. CONCLUSION AND FUTURE IMPROVEMENTS
% ============================================================================
\section{Conclusion and Future Improvements}

\subsection{Conclusion}

This project demonstrates the successful integration of deep learning edge detection with real-time hand gesture recognition to create an interactive, hands-free digital drawing and colouring application.
The custom-trained U-Net model produces high-quality contour maps from arbitrary images, and the multi-stage post-processing pipeline ensures that the resulting outlines are suitable for region-based colouring.
The gesture recognition system, built on geometric heuristics with debounce stabilisation and cursor smoothing, provides a responsive and intuitive interaction model.
The dual-mode architecture (HoloDraw and HoloPaint) showcases both traditional creative tools and AI-augmented workflows within a unified application.

\subsection{Future Improvements}

\begin{enumerate}[nosep]
    \item \textbf{GPU Acceleration:} Port the U-Net inference to GPU (CUDA/Metal) to enable real-time edge detection on video streams, enabling live ``edge overlay'' modes.

    \item \textbf{Advanced Edge Models:} Explore more recent architectures such as HED (Holistically-Nested Edge Detection) or BDCN (Bi-Directional Cascade Network) that may produce sharper, more complete boundaries.

    \item \textbf{Stroke Interpolation:} Implement cubic spline interpolation between cursor positions to produce smoother freehand strokes at high cursor velocities.

    \item \textbf{Multi-Hand Support:} Extend the gesture system to track both hands simultaneously, enabling bimanual interactions (e.g., one hand draws, the other adjusts brush size).

    \item \textbf{Collaborative Drawing:} Add network support to allow multiple users to draw on a shared canvas in real time.

    \item \textbf{Model Fine-Tuning:} Fine-tune the U-Net on domain-specific datasets (e.g., colouring book pages, cartoon line art) to improve edge quality for the target use case.
\end{enumerate}

% ============================================================================
% REFERENCES
% ============================================================================
\begin{thebibliography}{9}

\bibitem{bsds500}
P.~Arbelaez, M.~Maire, C.~Fowlkes, and J.~Malik,
``Contour detection and hierarchical image segmentation,''
\textit{IEEE Transactions on Pattern Analysis and Machine Intelligence},
vol.~33, no.~5, pp.~898--916, 2011.

\bibitem{ronneberger2015unet}
O.~Ronneberger, P.~Fischer, and T.~Brox,
``U-Net: Convolutional networks for biomedical image segmentation,''
in \textit{Proc. MICCAI}, Springer, 2015, pp.~234--241.

\bibitem{mediapipe}
C.~Lugaresi \textit{et al.},
``MediaPipe: A framework for building perception pipelines,''
\textit{arXiv preprint arXiv:1906.08172}, 2019.

\bibitem{opencv}
G.~Bradski,
``The OpenCV Library,''
\textit{Dr. Dobb's Journal of Software Tools}, 2000.

\end{thebibliography}

\end{document}
